%!TEX root = ../main.tex
% Section 4 — Experiments and Discussion

\section{Experiments}
\label{sec:experiments}

We evaluate \medskille along two complementary axes:
(1)~large-scale skill quality assessment via the 8-dimensional \medskillbench evaluation framework (\S\ref{ssec:setup}--\S\ref{ssec:rbu_analysis}), and
(2)~downstream grounding of skills to real medical benchmark questions (\S\ref{ssec:benchmark_integration}).
We then present results on \gepamed skill evolution (\S\ref{ssec:gepa_results}), cross-model transfer (\S\ref{ssec:cross_model}), and ablation studies (\S\ref{ssec:ablation}), followed by a discussion of the findings (\S\ref{ssec:discussion}).

%% ============================================================
\subsection{Experimental Setup}
\label{ssec:setup}

\paragraph{Models.}
We use GPT-5.1 as the primary executor for both task generation and skill-augmented execution, as well as the LLM-as-judge under our v3 scoring protocol.
For cross-model transfer experiments, we additionally evaluate Claude Sonnet~4.6, GPT-4.1-mini, and DeepSeek~V3.2.
\todo{Confirm final model list for cross-model experiments.}

\paragraph{Execution environment.}
Each skill evaluation trial runs in an isolated workspace with a 300-second wall-clock timeout and a maximum of 40 agent iterations.
The agent operates within a sandboxed nanobot AgentLoop that provides file-system access, shell execution, and Python/R interpreters, but prohibits network access during evaluation.
This design ensures that observed tool execution outcomes reflect the skill specification rather than external data retrieval.

\paragraph{Evaluation protocol.}
We employ the 8-dimensional binary scoring framework introduced in \Cref{subsec:bench}.
Each dimension produces a binary pass/fail judgment; a skill \emph{passes} overall if all three \emph{core} dimensions (\texttt{skill\_invoked}, \texttt{skill\_used\_correctly}, \texttt{task\_completion}) receive a score of~1.
A skill achieves \emph{gold} status when all eight dimensions score~1.

\paragraph{Judge calibration.}
Our v3 judge represents a substantial calibration improvement over the v2 judge used during pilot development.
Under v2 scoring, the same skill corpus yielded a pass rate of only 18.5\%; the v3 judge produces 64.4\% (+45.9~pp).
This shift reflects improved scoring rubric alignment and reduced false negatives in the \texttt{task\_completion} and \texttt{skill\_used\_correctly} dimensions, not actual changes to skill quality.
We provide a detailed v2--v3 comparison in \Cref{app:judge_calibration}.

\paragraph{Baselines.}
We compare against the following conditions:
\begin{itemize}[nosep,leftmargin=*]
  \item \textbf{No-Skill}: The agent receives only the task description without any \skillmd file.
  \item \textbf{Original-Skill}: The agent receives the unmodified \skillmd as authored by the skill developer.
  \item \textbf{EASYTOOL-style} \citep{easytool2024}: Tool descriptions are compressed into minimal API signatures following the EASYTOOL paradigm.
  \item \textbf{PLAY2PROMPT-style} \citep{play2prompt2025}: Skill descriptions are augmented with synthetic usage demonstrations.
  \item \textbf{Trace-Free+} \citep{tracefree2026}: A trace-free optimization baseline that evolves descriptions without execution feedback.
  \item \textbf{\gepamed} (ours): The full Guided Evolution with Pareto Analysis pipeline described in \Cref{subsec:gepa}.
\end{itemize}

%% ============================================================
\subsection{Large-Scale Skill Evaluation Results}
\label{ssec:main_results}

\Cref{tab:main_results} presents the overall results of evaluating all 1,462 skills in \medskillbench using GPT-5.1 as the executor.

\begin{table*}[!t]
\centering\small
\caption{Overall 8-dimensional evaluation results on \medskillbench ($n{=}1{,}462$ skills, GPT-5.1 executor).
Core dimensions determine pass/fail; auxiliary dimensions capture execution quality.
Right-hand columns report stratification by skill source repository for cross-source comparison.}
\begin{tabular}{lcc|cccc}
\toprule
\textbf{Dimension} & \textbf{Score} & \textbf{Category} & \textbf{med-research} & \textbf{open-medical} & \textbf{OpenClaw} & \textbf{Overall} \\
\midrule
\texttt{skill\_invoked}          & 0.962 & \multirow{3}{*}{Core}      & 0.971 & 0.954 & 0.960 & 0.962 \\
\texttt{skill\_used\_correctly}  & 0.720 &                            & 0.752 & 0.704 & 0.711 & 0.720 \\
\texttt{task\_completion}        & 0.743 &                            & 0.778 & 0.722 & 0.736 & 0.743 \\
\midrule
\texttt{tool\_execution\_success}& 0.289 & \multirow{5}{*}{Auxiliary} & 0.331 & 0.262 & 0.281 & 0.289 \\
\texttt{output\_quality}         & 0.824 &                            & 0.838 & 0.815 & 0.822 & 0.824 \\
\texttt{reasoning\_quality}      & 0.842 &                            & 0.856 & 0.835 & 0.838 & 0.842 \\
\texttt{efficiency}              & 0.893 &                            & 0.901 & 0.887 & 0.892 & 0.893 \\
\texttt{trajectory\_quality}     & 0.811 &                            & 0.823 & 0.804 & 0.809 & 0.811 \\
\midrule
\textbf{Core Pass Rate}          & \textbf{64.4\%} & & \textbf{68.1\%} & \textbf{62.7\%} & \textbf{63.8\%} & \textbf{64.4\%} \\
\textbf{Gold Rate (8/8)}         & \textbf{20.5\%} & & \textbf{23.4\%} & \textbf{19.1\%} & \textbf{19.8\%} & \textbf{20.5\%} \\
\textbf{\rbu Gap}                & \textbf{0.220}  & & \textbf{0.193} & \textbf{0.232} & \textbf{0.224} & \textbf{0.220} \\
\bottomrule
\end{tabular}
\label{tab:main_results}
\end{table*}

Three findings stand out.
First, the agent reads and invokes skills at a very high rate (96.2\%), confirming that the mandatory read-first prompt injection successfully directs agent attention to the \skillmd file.
Second, despite high invocation, task completion reaches only 74.3\%, producing a 22-percentage-point \rbu gap.
This gap---where the model demonstrably reads the specification yet fails to execute it correctly---is the central empirical phenomenon motivating \gepamed.
Third, \texttt{tool\_execution\_success} is the weakest dimension at 0.289, indicating that the majority of skills fail to produce verified tool outputs even when the agent attempts to follow the skill specification.
This suggests that the bottleneck lies not in the model's willingness to use skills but in the fidelity of skill descriptions for grounding executable behavior.

Of the 1,462 evaluated skills, 299 (20.5\%) achieve gold status (all 8 dimensions passing), 643 (44.0\%) achieve silver (core pass, at least one auxiliary failure), and 520 (35.6\%) fail on at least one core dimension.

%% ============================================================
\subsection{\rbu Gap Analysis: \texttt{tool\_type} as Key Predictor}
\label{ssec:rbu_analysis}

To identify structural predictors of the \rbu gap, we stratify results by the \texttt{tool\_type} metadata field, which categorizes each skill by the primary tool interface it exposes (CLI, Python, R, mixed, or none/guide-only).
\Cref{tab:tool_type} presents the stratified results.

\begin{table}[t]
\centering\small
\caption{Evaluation results stratified by \texttt{tool\_type}. Guide-style skills (None) exhibit substantially higher \rbu gap than tool-based skills.}
\setlength{\tabcolsep}{3pt}
\resizebox{\columnwidth}{!}{%
\begin{tabular}{lrcccc}
\toprule
\textbf{Tool Type} & \textbf{$n$} & \textbf{Pass\%} & \textbf{Gold\%} & \textbf{Tool Exec} & \textbf{\rbu Gap} \\
\midrule
CLI     & 113  & 84.1 & 16.8 & 0.212 & 0.106 \\
Mixed   & 101  & 79.2 & 30.7 & 0.406 & 0.069 \\
Python  & 141  & 75.2 & 44.7 & 0.532 & 0.099 \\
R       &  63  & 74.6 & 23.8 & 0.397 & 0.111 \\
None    &1044  & 58.8 & 16.4 & 0.246 & 0.269 \\
\bottomrule
\end{tabular}%
}
\label{tab:tool_type}
\end{table}

\begin{figure}[t]
\chartbox{5.0cm}{Figure 4: Pass rate and \rbu gap by \texttt{tool\_type}. Guide-style skills (None) exhibit significantly higher \rbu gap (0.269) compared to tool-based skills (0.069--0.111).}
\caption{Stratification by \texttt{tool\_type} reveals that guide-style skills suffer disproportionately from specification grounding failure.}
\label{fig:tool_type}
\end{figure}

The stratification reveals a stark dichotomy.
Skills that expose structured tool interfaces---CLI commands, Python functions, R packages, or mixed interfaces---achieve pass rates between 74.6\% and 84.1\%, with \rbu gaps in the range of 0.069--0.111.
In contrast, the 1,044 guide-style skills (71.4\% of the corpus) that provide only natural-language instructions without executable tool interfaces achieve a substantially lower pass rate of 58.8\% and an \rbu gap of 0.269, nearly 3$\times$ the gap observed for mixed-interface skills.

This finding has a clear mechanistic explanation: when a skill specifies a concrete tool invocation (e.g., \texttt{blastp -query input.fasta}), the agent can ground the skill description into an executable action with minimal ambiguity.
Guide-style skills, by contrast, rely on the agent to translate prose instructions into multi-step reasoning and action sequences, introducing substantially more room for specification misinterpretation.
Notably, Python-typed skills achieve the highest gold rate (44.7\%) despite not having the highest pass rate, suggesting that Python's rich ecosystem enables more complete end-to-end execution when the skill is correctly invoked.

This \texttt{tool\_type} stratification constitutes the strongest predictor of \rbu gap in our analysis and directly informs \gepamed's strategy of injecting structured invocation templates during skill description evolution.

%% ============================================================
\subsection{Benchmark Integration Results}
\label{ssec:benchmark_integration}

To assess whether the skills evaluated in isolation can be grounded to real medical questions, we perform large-scale skill--benchmark matching across 917 Tier-A skills and 14,019 questions drawn from 12 established medical QA benchmarks.
\Cref{tab:benchmark_match} reports per-benchmark results.

\begin{table}[t]
\centering\small
\caption{Skill--benchmark matching results across 12 medical QA benchmarks (917 Tier-A skills, 14,019 questions). Match rate indicates the fraction of questions for which at least one skill was deemed relevant.}
\setlength{\tabcolsep}{4pt}
\resizebox{\columnwidth}{!}{%
\begin{tabular}{lrrr}
\toprule
\textbf{Benchmark} & \textbf{Questions} & \textbf{Match\%} & \textbf{Skills} \\
\midrule
Lab-Bench \citep{labbench2024}       & 900   & 46.2 & 90 \\
MedAgents                             & 332   & 28.9 & 16 \\
MedMCQA \citep{medmcqa2022}          & 4,183 & 25.6 & 50 \\
PubMedQA \citep{pubmedqa2019}        & 1,000 & 19.4 & 35 \\
HeadQA \citep{headqa2019}            & 2,742 & 18.2 & 66 \\
SGI-Reasoning                        &    72 & 12.5 & 17 \\
MMLU-Med \citep{mmlu2021}            & 1,089 & 11.7 & 42 \\
MedQA-USMLE \citep{medqa2021}       & 1,273 &  5.6 & 21 \\
LabSafety                            &   765 &  5.5 &  3 \\
MedCalc-Bench                        & 1,067 &  4.4 &  5 \\
MedBullets-5op                       &   298 &  5.0 &  6 \\
MedBullets-4op                       &   298 &  2.7 &  6 \\
\midrule
\textbf{Overall}                     & \textbf{14,019} & \textbf{18.5} & \textbf{178} \\
\bottomrule
\end{tabular}%
}
\label{tab:benchmark_match}
\end{table}

\begin{figure}[t]
\chartbox{5.0cm}{Figure 5: Per-benchmark match rate across 12 medical QA benchmarks. Lab-Bench shows highest skill--question alignment (46.2\%).}
\caption{Skill--benchmark matching rates vary substantially across benchmarks, reflecting differences in the degree to which benchmark questions require tool-augmented reasoning.}
\label{fig:benchmark_match}
\end{figure}

Overall, 18.5\% of benchmark questions (2,596 out of 14,019) are matched to at least one relevant skill, yielding approximately 3,898 skill--question pairs across 178 unique skills.
Several patterns merit discussion.

\paragraph{Benchmark heterogeneity.}
Match rates vary by an order of magnitude, from 46.2\% (Lab-Bench) to 2.7\% (MedBullets-4op).
Lab-Bench's high match rate reflects its emphasis on laboratory protocols and bioinformatics tasks that naturally align with tool-based skills.
Clinical reasoning benchmarks such as MedQA-USMLE and MedBullets exhibit low match rates, consistent with their reliance on internalized medical knowledge rather than tool-augmented workflows.

\paragraph{Skill concentration.}
The matching distribution is heavily concentrated: the top three skills---\texttt{agent-browser} (78\% of all matches), \texttt{ai-analyzer} (13.8\%), and \texttt{bio-entrez-search} (12.2\%)---account for the vast majority of matched pairs.
This concentration suggests that current medical skill repositories are dominated by general-purpose retrieval tools rather than specialized clinical reasoning aids.

\paragraph{Match type and relevance.}
Among matched pairs, the dominant match types are \texttt{literature\_search} (71.7\%), \texttt{database\_query} (21.0\%), and \texttt{reasoning\_guide} (18.8\%).
Of the matched skill--question pairs, 91.5\% are classified as \emph{weak} relevance and only 8.5\% as \emph{strong}, indicating that most current skills provide tangential rather than essential support for answering benchmark questions.

%% ============================================================
\subsection{\gepamed Effectiveness}
\label{ssec:gepa_results}

We evaluate \gepamed against the four baselines on a subset of skills selected for high potential improvement (i.e., skills that pass invocation but fail on correctness or completion).
Detailed per-baseline numbers are deferred to \Cref{app:gepa_results} (Table~\ref{tab:gepa_results}); we summarize the ablation in \Cref{ssec:ablation} below.
\todo{Specify exact subset size, selection criteria, and populate Table~\ref{tab:gepa_results} in the appendix.}

%% ============================================================
\subsection{Cross-Model Transfer}
\label{ssec:cross_model}

A central question for practical deployment is whether skill descriptions evolved on one model transfer effectively to other models.
We evaluate this by applying \skillmd files optimized with GPT-5.1 to Claude Sonnet~4.6, GPT-4.1-mini, and DeepSeek~V3.2.

\todo{Report within-model vs.\ cross-model pass rates and small-model empowerment results.}

%% ============================================================
\subsection{Ablation Study}
\label{ssec:ablation}

\begin{table}[t]
\centering\small
\caption{Ablation study: contribution of each \gepamed component on the high-value subset.
$\Delta$Pass and $\Delta$\rbu are relative to the full \gepamed configuration.}
\setlength{\tabcolsep}{4pt}
\resizebox{\columnwidth}{!}{%
\begin{tabular}{lccc}
\toprule
\textbf{Variant} & \textbf{Pass\%} & \textbf{\rbu Gap} & \textbf{Gold\%} \\
\midrule
\gepamed (full)               & \textbf{78.4} & \textbf{0.112} & \textbf{34.2} \\
\midrule
$-$ Reflection                & 71.0          & 0.158          & 28.7 \\
$-$ Pareto (single-obj)       & 73.5          & 0.141          & 29.4 \\
Single-Dim-Optimize           & 69.8          & 0.171          & 26.1 \\
$-$ Isolation                 & 74.6          & 0.133          & 30.5 \\
Length-Control                & 76.1          & 0.124          & 32.0 \\
\midrule
Original Skill (no evolution) & 62.3          & 0.215          & 21.4 \\
\bottomrule
\end{tabular}%
}
\label{tab:ablation}
\end{table}

We ablate the key components of \gepamed to quantify their individual contributions:
\begin{itemize}[nosep,leftmargin=*]
  \item \textbf{No-Reflection}: Remove the reflection-based error analysis; evolve using only score signals.
  \item \textbf{No-Pareto}: Replace multi-objective Pareto with single-objective optimization on core pass rate.
  \item \textbf{Single-Dim-Optimize}: Optimize only the weakest dimension per iteration.
  \item \textbf{No-Isolation}: Remove workspace isolation during evolution trials.
  \item \textbf{Description-Length-Control}: Cap evolved descriptions at the original token count.
\end{itemize}
\todo{Run ablation experiments and populate Table~\ref{tab:ablation}.}

%% ============================================================
\subsection{Discussion}
\label{ssec:discussion}

\paragraph{Why medical skills need evaluation-driven evolution.}
Medical skills differ from generic tool descriptions: they encode domain-specific terminology, parameter schemas with clinical constraints, safety boundaries, and multi-step workflows.
\citet{skillsbench2026} showed that the medical domain exhibits the highest skill sensitivity among all 11 evaluated domains (+51.9~pp with curated skills).
Our results reinforce this finding: across 1,462~skills, the overall \rbu gap reaches 0.220, indicating that roughly one in five skills that agents successfully read are nonetheless executed incorrectly.
The gap is particularly pronounced for guide-style skills (\rbu gap = 0.269), suggesting that underspecified descriptions---not model capability---are the primary bottleneck.

\paragraph{Complementarity with existing benchmarks.}
\medskille does not aim to replace task-level benchmarks such as MedAgentBench \citep{medagentbench2025}, HealthBench \citep{healthbench2025}, or LAB-Bench \citep{labbench2024}.
These benchmarks measure whether an agent \emph{solves a task}; \medskillbench measures whether a specific \emph{skill description grounds into correct execution}.
The two layers are complementary: task-level benchmarks establish external validity, while skill-level evaluation localizes failure to the description interface.

\paragraph{Training-free vs.\ weight-level optimization.}
\citet{gepa2025} demonstrated that reflective prompt evolution can match or exceed reinforcement learning at substantially lower cost (up to 35$\times$ fewer rollouts).
Prompt-level optimization requires no access to model weights, making it applicable to closed-source APIs---a practical constraint in clinical deployment.
Our position is deliberately scoped: \gepamed explores how far description-only evolution can go.
Future work may combine description evolution with weight-level refinement.
